\begin{table}
\caption{Marginal CV effect ($\tau$) and voluntary-side baseline by sex $\times$ SES cell (age-18 cutoff, pp).}
\label{tab:triple_cells_T18}
\begin{tabular}{lllrrrr}
\toprule
SES measure & Sex & Deprivation & $\tau$ (pp) & SE ($\tau$) & Baseline (pp) & SE (baseline) \\
\midrule
NBI & Men & Low (p10) & 26.58 & 0.87 & 47.26 & 0.73 \\
NBI & Men & High (p90) & 17.92 & 0.98 & 48.83 & 0.82 \\
NBI & Women & Low (p10) & 26.53 & 0.85 & 50.20 & 0.69 \\
NBI & Women & High (p90) & 14.61 & 0.98 & 56.74 & 0.82 \\
PCA index & Men & Low (p10) & 28.84 & 0.97 & 47.91 & 0.82 \\
PCA index & Men & High (p90) & 16.20 & 0.97 & 48.06 & 0.81 \\
PCA index & Women & Low (p10) & 29.57 & 0.95 & 49.04 & 0.78 \\
PCA index & Women & High (p90) & 12.39 & 0.98 & 57.42 & 0.86 \\
\bottomrule
\end{tabular}
\par\vspace{0.6em}
{\footnotesize\begin{minipage}{\linewidth}\emph{Note:} Cells are marginal effects from the continuous triple-interaction model (Section~\ref{sec:strat_hetero}, Equation~\ref{eq:triple}), evaluated at the 10th and 90th percentiles of the SES measure within the pooled MSE-optimal window ($\approx$66 days). They are not the stratified per-cell RD estimates of Table~\ref{tab:rdd_results}, each of which is fit at its own bandwidth, and the two need not coincide. Standard errors clustered by electoral circuit (CRV1).\end{minipage}}
\end{table}
